%!TEX root = ../../../main.tex

\chapter{Einführung}
\label{chap:einfuehrung}

\section{Grundgedanke und Motivation}
\label{sec:motivation}

Fast jeder Haushalt hat einen Ordner mit Rechnungen und Verträgen für Strom, Wasser, Handy, Versicherungen oder Arztbesuche, teils auf Papier, teils als PDF. Hineingeschaut wird erst, wenn man etwas braucht, und dann wird es oft mühsam oder teuer. Wer die Kündigungsfrist eines Vertrags verpasst, zahlt häufig ein ganzes weiteres Jahr. Ob eine erhöhte Abschlagszahlung zum Vertrag passt, merkt man nur, wenn man Rechnung und Vertrag nebeneinanderlegt.

Naheliegend wäre ein Chatbot in der Cloud. Gerade diese Unterlagen gehören aber zu den privatesten Daten überhaupt. Darin stehen Name, Adresse und Geburtsdatum, die IBAN, Kunden- und Vertragsnummern, auf Arztrechnungen sogar Behandlungen, und der Stromverbrauch verrät ziemlich genau, wann jemand zu Hause ist. Meine Grundidee ist deshalb ein Agent, der komplett auf meinem Laptop läuft, ohne Internet und Grafikkarte, mit 16\,GB Arbeitsspeicher.

Auf so einem Rechner laufen nur kleine Sprachmodelle mit 4 bis 8 Milliarden Parametern. In meinen ersten Versuchen haben sie Texte erstaunlich gut gelesen, beim Rechnen und bei Datumsangaben aber oft danebengelegen. Bei Geld und Fristen ist eine falsche Antwort, die sicher klingt, schlimmer als gar keine. In meiner Architektur liest und formuliert deshalb das Modell, und alles, was sich ausrechnen oder nachprüfen lässt, übernimmt Python-Code, etwa Fristen, Filter und die Prüfung, ob eine Zahl im Dokument steht.

\section{Projektidee und Anwendungsfall}
\label{sec:usecase}

So ist LocalLLM entstanden, eine Streamlit-App für den privaten Dokumentenordner. Man lädt PDF- oder Word-Dateien hoch, und das System liest jedes Dokument einmal ein. Über einen Hash merkt es sich, welche Dateien es schon kennt, sodass eine neue Rechnung allein verarbeitet wird (auf meinem Laptop etwa zweieinhalb Minuten). Danach hilft die App bei drei Dingen, die im Alltag Zeit kosten. Das erste ist das \emph{Finden} einer einzelnen Angabe, etwa des Arbeitspreises im Stromvertrag oder des Zahlungsziels einer Arztrechnung. Das zweite ist das \emph{Prüfen}, etwa ob die Wasserrechnung mit der vereinbarten Vorauszahlung rechnet. Das dritte ist der \emph{Überblick}, etwa welche Verträge sich bis Jahresende kündigen lassen. Dazu gibt es eine Vertragsübersicht, die für jeden Vertrag den letzten Kündigungstag ausrechnet und per Ampel zeigt, wo es in den nächsten 30 oder 90 Tagen eng wird.

Wichtig war mir, dass man jede Antwort nachprüfen kann. Ein Klick auf die Quelle zeigt die Seite des Originaldokuments mit dem Wert gelb markiert, und ohne passendes Dokument soll der Agent das sagen statt zu raten. Meine echten Unterlagen konnte ich hier nicht verwenden. Ich habe deshalb 16 Dokumente mit realistischem Layout aus festen Sollwerten erzeugt, zehn Rechnungen als PDF und sechs Verträge (drei DOCX, drei PDF). So kenne ich zu jeder Testfrage Antwort und Dokument. Absichtlich schwierig sind etwa Hoch- und Niedertarif in Nachbarzeilen, Netto- und Bruttospalte nebeneinander und Stromwerte, die nur in einem Diagramm stehen. Die Fragen an solche Dokumente lassen sich in fünf Typen einteilen (Tabelle~\ref{tab:fragetypen}), nach denen sich die Architektur richtet.

\begin{table}[tb]
    \centering\footnotesize\setstretch{1.05}
    \begin{tabularx}{\textwidth}{@{}l Y Y@{}}
        \toprule
        \textbf{Typ} & \textbf{Beispiel} & \textbf{Was der Agent dafür braucht} \\
        \midrule
        a) Wert aus Tabelle & USt.\ der LuminaStyle-Rechnung? & richtige Zeile mit Spaltenkopf \\
        b) Textstelle & Mindestlaufzeit Handyvertrag? & richtigen Abschnitt im Vertrag \\
        c) Zwei Dokumente & Vorauszahlung wie im Vertrag? & Rechnung \emph{und} Vertrag \\
        d) Liste & Welche Verträge verlängern sich? & alle Verträge, berechnete Fristen \\
        e) ohne Beleg & Frist meines Vodafone-Vertrags? & erkennen, dass nichts passt \\
        \bottomrule
    \end{tabularx}
    \caption{Fragetypen und was sie vom Agenten verlangen}
    \label{tab:fragetypen}
\end{table}

Gebaut habe ich zwei Agenten auf denselben Daten. Der \emph{Orchestrator} arbeitet jede Frage in einem festen Ablauf ab, den der Code vorgibt, ein ReAct-Agent entscheidet dagegen selbst, welches Werkzeug er wann benutzt. Ich wollte wissen, ob das kleinen Modellen hilft.
